\section{Introduction}
\label{sec:introduction}

A scientific agent that correctly predicts its own failure has solved only
half the reliability problem. The other half is deciding what to do about
it: whether to buy additional evidence, which verification check to run,
whether recovery is still possible at acceptable cost, and when to escalate
or abstain. A calibrated warning that arrives late in the trajectory,
after the recovery window has closed, produces no practical benefit. A
warning that triggers unnecessary verification on every clean task wastes
compute and risks false interventions. The scientific question is not
whether failure can be detected—it can—but whether a policy can use
calibrated detection results to choose cost-effective verification and
recovery actions that improve outcomes without degrading performance on
tasks that would have succeeded without intervention.

Recent work on agent reliability provides important foundations but stops
short of this question. AgentRx localizes critical steps and attributes
root causes from execution traces~\citep{agentrx2026}. Aegis pairs a
failure taxonomy with intervention suggestions~\citep{aegis2025}. CAPF
converts zero-reward failures into verifier-guided repair
trajectories~\citep{capf2026}. Afterburner compares SFT, DPO, and GRPO
under execution feedback~\citep{afterburner2025}. These systems demonstrate
that diagnosis, repair, and execution-feedback training are individually
feasible. However, none provides a factorized decision framework in which a
budget-aware policy selects verification or recovery actions from calibrated
beliefs, the outcomes are independently adjudicated, and the entire
pipeline—belief, decision, verification, recovery, and training—is tested
on frozen unseen fault compositions with a clean-task safety gate.

This paper makes two separately identifiable claims. First, a calibrated,
cost-sensitive verification policy improves independently verified
scientific success at matched verification budget relative to fixed, random,
and confidence-threshold policies, evaluated in a randomized online study
with task-level blocking (H3). Second, process-supervised training on
independently adjudicated recovery transitions improves unseen-fault
recovery and selective risk at matched cost beyond recovery SFT and
outcome-only optimization, without meaningful regression on clean tasks
(H4). Both claims are evaluated under the frozen trace contract, fault
model, and held-out splits of the companion benchmark paper~\citep{reliability_benchmark}, which
supplies calibrated failure hazard and cause beliefs as fixed inputs.
The benchmark contract is consumed without modification; R2 may add
intervention fields through a new schema version but may not relabel
R1's test outcomes.

The paper structure follows the dependency. After positioning against
related work (Section~\ref{sec:related_work}), we formalize the sequential
decision problem over belief states, actions, costs, and postconditions
(Section~\ref{sec:problem}). The controller and adjudication method is
presented in Section~\ref{sec:method}. The verified-intervention dataset
and randomized study design are specified in Section~\ref{sec:dataset}.
The full experimental matrix—six controller arms, six post-training arms,
matched budgets, and evaluation sets—is detailed in
Section~\ref{sec:experiments}. Results appear in
Section~\ref{sec:results} and are discussed in
Section~\ref{sec:discussion}.
